\begin{textsample}{NEG Dim 8 – global\_north\_2024\_02 – Score -1.00 – global\_north\_2024\_02/105918402.txt}  \label{ex:f8_neg_438}
The international disaster relief charity ShelterBox is set to support thousands of people displaced in Gaza with emergency shelter aid and other essential items.
In partnership with Medical Aid for Palestinians ( MAP ), the charity will be providing tarpaulins, rope, and other items so that people can make temporary repairs to damaged buildings that will help keep them watertight and protect them from the weather.
The humanitarian need is growing, with chilling winter temperatures putting the lives of 1.9 million displaced people -- more than 85 \% of the population in Gaza - at risk.
People in Gaza were already reliant on humanitarian aid and the current volume of aid crossing into Gaza is fraction of what ' s needed.
ShelterBox ' s Regional Director for the Middle East and North Africa ( MENA ), Haroon Altaf says :"Up to 1.9 million people were uprooted from their homes over a \textbf{two-month} period -- the scale of the crisis is huge.
Let us not forget that behind every number are individuals -- children, women, men -- who have a name, a family, a @ @ @ @ @ @ @ @ @ @ hunger, and the cold."Winter makes things harder in every emergency.
It ' s cold and wet and people don't have what they need to protect themselves or their families from the harsh winter weather."As well as emergency shelter, ShelterBox and MAP will be providing blankets, mattresses, pillows, and floor mats to help people stay warmer and save lives.
The aid package will include washing sets, water carriers, kitchen sets, and items like nappies, toothbrushes, sanitary items, and body/hair wash.
Vital items for people with no belongings and unable to return home, or people who have had their homes damaged.
Are you getting our free newsletter?
ShelterBox has launched an urgent fundraising appeal to provide emergency shelter and other essential items to people left with nowhere to live in Gaza, and other disasters around the world.
More than 3.1 million people are in need across Gaza and the West Bank.
Almost 1.4 million people are sheltering in designated shelters like @ @ @ @ @ @ @ @ @ @ shelters or out in the open in the street.
There is severe overcrowding.
Disease is spreading and that ' s placing extra strain on the already overwhelmed health system.
Getting aid into countries and on to where it ' s needed during times of conflict can be difficult, especially with damage to infrastructure and supply routes.
Access to Gaza is limited, so to provide support to people displaced, ShelterBox is partnering with MAP - a medical organisation that was already working in Gaza before the war started.
Working together with them, means ShelterBox can get emergency shelter aid and other essential aid items into Gaza, and to the people who need it, more quickly.
ShelterBox ' s Programme Manager for Gaza, Jonty Ellaby adds :"Getting aid into Gaza is challenging.
It is slow and unpredictable, with border crossings often closed at short notice."The situation is complex, but we intend that our aid will reach people in the coming weeks via the Rafah border crossing between Egypt and Gaza. @ @ @ @ @ @ @ @ @ @ are prioritised for entry, with aid moved into Gaza from Egypt and Jordan via Rafah or the Kerem Shalom border crossing between Israel and Gaza."ShelterBox has experience getting aid into Gaza and worked there in 2004, 2008, and 2015.
In 2015, ShelterBox distributed emergency shelter aid, thermal blankets, and winter clothing to people in Gaza.
Supporting ShelterBox will help people affected displaced in Gaza and other disasters around the world.
The charity believes it can have the most impact for communities by staying flexible.
It allows ShelterBox to restock its warehouses ready for the next disaster and provide support to people affected by extreme weather, conflict, and climate crises around the world.
Did you know Scoop has an Ethical Paywall?
If you ' re using Scoop for work, your organisation needs to pay a small license fee with Scoop Pro.
We think that ' s fair, because your organisation is benefiting from using our news resources.
In return, we ' ll also give your team access to @ @ @ @ @ @ @ @ @ @, because public access to news is important!

% matched lemmas: two-month
\end{textsample}
